\chapter{Couche 6 — La qualité}

La plateforme marche. Reste à garantir qu'elle marchera encore après la
prochaine modification --- et qu'elle ne livre pas de faille connue. Treize
briques, dont onze tournent à chaque envoi de code.

% -----------------------------------------------------------------------------
\outil{pytest}{brique 34 · les tests}

\quoi
Le cadre de test de Python. Une fonction dont le nom commence par \cmd{test\_}
est un test ; une assertion fausse le fait échouer.

\pourquoi
L'alternative dans la bibliothèque standard est \cmd{unittest}, qui impose
d'écrire des classes et des méthodes à noms imposés. pytest accepte des
fonctions simples, et surtout il apporte les \emph{fixtures} : des morceaux de
préparation réutilisables, dont l'un est central ici.

\role
168 tests, en vingt secondes, répartis en six fichiers. Deux mécanismes font
tout le travail.

\textbf{La base vierge à chaque test.} Une fixture marquée \cmd{autouse=True}
efface et recrée la base avant chaque test, sans qu'on ait à y penser.

\begin{piege}{Des tests qui dépendent de leur ordre}
Si un test laisse des données derrière lui, le suivant peut réussir pour la
mauvaise raison --- ou échouer quand on le lance seul. Le symptôme : « ça passe
en entier mais pas tout seul », ou l'inverse. C'est le genre de bug qui fait
abandonner une suite de tests.
\end{piege}

\textbf{L'espion sur Docker.} Les tests ne lancent aucun conteneur :
\cmd{run\_docker} est remplacé par une fonction qui note ce qu'on lui demande.
On vérifie ainsi que la commande imposée par le sujet part bien telle quelle,
sur une machine qui n'a même pas Docker.

\apport
Les tests tournent partout, en vingt secondes. C'est ce qui permet de les lancer
à chaque envoi plutôt qu'une fois par mois.

% -----------------------------------------------------------------------------
\outil{ruff}{brique 35 · l'analyse de style et d'erreurs}

\quoi
Un analyseur qui lit le code sans l'exécuter et signale ce qui est douteux :
import jamais utilisé, variable inconnue, tournure obsolète, motif dangereux.

\pourquoi
Les outils historiques sont \cmd{flake8} pour le style et \cmd{pylint} pour
l'analyse, plus \cmd{isort} pour les imports. Ruff fait les trois, est écrit en
Rust, et tourne en moins d'une seconde sur tout le projet --- contre plusieurs
dizaines pour pylint. Cette vitesse change l'usage : on le lance avant chaque
commit au lieu de l'éviter.

\role
Configuré dans \fichier{pyproject.toml}, avec six familles de règles
sélectionnées. Une seule est désactivée --- la longueur de ligne --- parce que
les commentaires pédagogiques du projet dépassent parfois.

Il vérifie aussi \cmd{target-version = "py310"} : il refuse les tournures que le
Python de la machine cible ne comprendrait pas.

\apport
Le code reste lisible et cohérent sans discussion de style, et les fautes
d'inattention sont attrapées avant l'exécution.

% -----------------------------------------------------------------------------
\outil{bandit}{brique 36 · l'analyse de sécurité du code (SAST)}

\quoi
\emph{Static Application Security Testing} : il cherche dans le code source des
motifs connus pour être dangereux.

\pourquoi
Un relecteur humain oublie. bandit connaît une liste de pièges classiques et ne
se fatigue pas : concaténation SQL, \cmd{shell=True}, secret en dur, algorithme
de hachage obsolète, fichier temporaire prévisible.

\role
Lancé avec \cmd{--severity-level low --confidence-level low} : on veut tout
voir. Le projet ne désactive \textbf{aucune règle globalement} ; les rares cas
légitimes sont annotés ligne à ligne avec leur justification --- par exemple
\cmd{\# nosec B404} sur l'import de \cmd{subprocess}, puisque le sujet impose de
piloter Docker par sa ligne de commande.

\begin{piege}{Le témoin de vie dans \texttt{/tmp}}
Le Faucheur horodate un fichier à chaque cycle, pour que la sonde du conteneur
sache que sa boucle tourne. Nous l'avions mis dans \cmd{/tmp} ; bandit a
signalé la règle B108, et il avait raison.

Le Faucheur tourne aussi en service systemd sur le contrôleur, où \cmd{/tmp} est
partagé avec tous les comptes locaux : un lien symbolique y suffirait à nous
faire horodater un fichier qui ne nous appartient pas. Le témoin vit désormais
dans le dossier de l'application.
\end{piege}

\apport
Une famille entière d'erreurs est attrapée automatiquement, y compris celles
qu'on vient d'introduire sans y penser.

% -----------------------------------------------------------------------------
\outil{pip-audit}{brique 37 · l'audit des dépendances (SCA)}

\quoi
\emph{Software Composition Analysis} : il confronte la liste des dépendances
épinglées aux bases de vulnérabilités publiques.

\pourquoi
On écrit soi-même quelques milliers de lignes, mais on en importe des
centaines de milliers. Une faille dans Flask ou dans Werkzeug est une faille du
projet, et on ne l'apprend pas en relisant son propre code.

L'alternative est \cmd{safety}, dont la base complète est payante. pip-audit est
publié par l'autorité de packaging Python et utilise les bases publiques.

\role
Lancé sur \fichier{requirements.txt} avec \cmd{--strict}. C'est la brique qui
rend l'épinglage des versions (brique 3) exploitable : on ne peut auditer que ce
qu'on connaît exactement.

\begin{retenir}
Ce travail trouve des failles \textbf{sans qu'on ait rien changé}. Werkzeug
3.1.8 est devenu vulnérable (CVE-2026-102598) entre deux exécutions de la
chaîne : notre code n'avait pas bougé, le monde si. Sans ce contrôle
automatique, nous l'aurions appris trop tard.
\end{retenir}

\apport
Les failles des dépendances sont vues le jour de leur publication, pas le jour
de l'incident.

% -----------------------------------------------------------------------------
\outil{gitleaks}{brique 38 · la recherche de secrets}

\quoi
Un outil qui parcourt le dépôt \textbf{et tout son historique} à la recherche de
clés, jetons et mots de passe.

\pourquoi
C'est la brique qui corrige une illusion : effacer un secret dans un commit
ultérieur ne l'efface pas du dépôt. Il reste dans l'historique, consultable par
quiconque clone.

\role
Lancé sur l'historique complet à chaque envoi. C'est ce qui permet au Vault
d'être versionné sans crainte : s'il était un jour commis en clair, la chaîne le
verrait immédiatement.

\apport
Une erreur humaine --- commettre un fichier de configuration contenant un mot de
passe --- est détectée avant qu'elle ne devienne publique.

% -----------------------------------------------------------------------------
\outil{nuclei et dast.sh}{brique 39 · l'analyse dynamique (DAST)}

\quoi
\emph{Dynamic Application Security Testing} : on attaque l'application
\textbf{en fonctionnement}, par le réseau, comme le ferait un attaquant.

\pourquoi
bandit lit le code ; il ne peut pas savoir si l'en-tête de sécurité est
réellement envoyé, ni si une page protégée l'est vraiment une fois déployée.
Seul un test contre le service qui tourne le dit.

\cmd{nuclei} est un scanner générique avec des milliers de signatures.
\fichier{tests/dast.sh} est un script maison qui vérifie ce que nuclei ne
connaît pas : les règles propres à ce projet.

\role
La chaîne démarre l'application avec Gunicorn, puis lance les deux. Dix-huit
contrôles : en-têtes de sécurité, pages protégées inaccessibles sans session,
jetons CSRF exigés, attributs des cookies.

\apport
On vérifie le comportement réel, pas l'intention du code.

% -----------------------------------------------------------------------------
\outil{hadolint et ansible-lint}{brique 40 · l'analyse de l'infrastructure}

\quoi
Deux analyseurs spécialisés : \cmd{hadolint} pour les Dockerfiles,
\cmd{ansible-lint} pour les rôles Ansible.

\pourquoi
Le code d'infrastructure est du code, avec ses propres pièges. Un Dockerfile
peut figer un mot de passe dans une couche, installer une version non épinglée,
ou tourner en root sans raison. Un rôle Ansible peut être non idempotent sans
que cela se voie.

\role
\cmd{hadolint} analyse les quatre Dockerfiles. C'est lui qui a trouvé le mot de
passe par défaut de la brique 16.

\cmd{ansible-lint} tourne au profil \cmd{production}, le plus strict. Il vérifie
notamment que les tâches ont un nom, que les modules sont qualifiés
(\cmd{ansible.builtin.*}), et signale les usages non idempotents.

\begin{piege}{\texttt{command} et \texttt{shell} ne sont jamais idempotents}
Ces deux modules exécutent ce qu'on leur donne sans rien vérifier : ils sont
donc \emph{toujours} signalés \cmd{changed}. Il faut leur apprendre à se taire
avec \cmd{creates:}, \cmd{removes:} ou \cmd{changed\_when:}.
\end{piege}

\apport
L'infrastructure est relue avec la même exigence que le code applicatif.

% -----------------------------------------------------------------------------
\outil{Trivy}{brique 41 · le scan des images}

\quoi
Un scanner qui inventorie les paquets d'une image Docker et cherche leurs
vulnérabilités connues.

\pourquoi
pip-audit couvre les dépendances Python. Mais une image contient aussi un
système de base --- Debian, Ubuntu, Alpine --- avec ses centaines de paquets.
Une faille OpenSSL dans l'image de base est une faille livrée.

\role
La chaîne inventorie \emph{toutes} les failles sans bloquer, puis relance Trivy
avec un verrou : aucune faille \cmd{HIGH} ou \cmd{CRITICAL}
\textbf{corrigeable} ne passe.

\begin{retenir}
Le mot « corrigeable » fait tout. L'option \cmd{--ignore-unfixed} écarte les
failles sans correctif disponible : les refuser ne protégerait de rien et
bloquerait tout, ce qui finirait par faire désactiver le contrôle.
\end{retenir}

Ce verrou a bloqué nos propres images : la base \cmd{docker:28-cli} traînait
163 failles corrigeables liées à sa version de Go. Il a fallu passer à
\cmd{docker:29-cli} --- deux failles au lieu de 163 --- et ajouter un
\cmd{apt-get upgrade} pour OpenSSL.

\apport
Aucune image vulnérable ne part au registre. Et la preuve que le contrôle sert :
il nous a refusé les nôtres.

% -----------------------------------------------------------------------------
\outil{Syft}{brique 42 · l'inventaire de ce qu'on livre (SBOM)}

\quoi
\emph{Software Bill of Materials} : la liste exhaustive des composants d'une
image livrée, avec leurs versions.

\pourquoi
Quand une faille majeure est annoncée, la première question est : « sommes-nous
touchés ? ». Sans inventaire, il faut reconstruire l'image et chercher. Avec un
SBOM archivé, on lit un fichier.

C'est aussi une exigence réglementaire croissante pour tout logiciel livré.

\role
Généré à chaque publication et conservé en artefact de la chaîne --- environ
340 Ko décrivant tout ce que contiennent les images.

\apport
Le projet peut répondre, des mois plus tard, à une question sur une faille
publiée après sa livraison.

% -----------------------------------------------------------------------------
\outil{GitHub Actions}{brique 43 · l'orchestration de tout cela}

\quoi
Le service qui exécute automatiquement une suite de travaux à chaque envoi de
code, sur des machines neuves.

\pourquoi
Les onze contrôles précédents ne valent que s'ils tournent. Lancés à la main,
ils sont oubliés le jour où l'on est pressé --- c'est-à-dire le jour où l'on en
a le plus besoin.

Les alternatives --- GitLab CI, Jenkins --- auraient demandé un serveur à
administrer. GitHub Actions est intégré au dépôt, et son exécution sur une
machine neuve à chaque fois garantit qu'on ne dépend pas d'un reste de
configuration locale.

\role
Onze travaux dans \fichier{.github/workflows/ci.yml}, 601 lignes. Et surtout une
\textbf{barrière} : la publication et le déploiement ne s'exécutent que si les
neuf précédents sont au vert, et seulement sur la branche principale.

\begin{retenir}
Sans dépendance déclarée (\cmd{needs:}), GitHub Actions lance tous les travaux
en parallèle --- et publierait l'image pendant que les tests échouent. C'est la
dépendance explicite qui transforme une liste de contrôles en véritable
barrière.
\end{retenir}

\begin{piege}{La chaîne vérifie sa propre syntaxe}
Une erreur de syntaxe shell dans un bloc \cmd{run:} ne se voit qu'au moment où
l'étape s'exécute, parfois après dix minutes. Trois exécutions ont échoué pour
deux guillemets déséquilibrées.

Un contrôle a été ajouté : il extrait chaque bloc \cmd{run:}, neutralise les
expressions de modèle, et passe le tout à \cmd{bash -n} avant que quoi que ce
soit ne démarre.
\end{piege}

\apport
Onze contrôles qui tournent vraiment, à chaque fois, sur une machine neuve.

% -----------------------------------------------------------------------------
\outil{ghcr.io et docker compose}{brique 44 · livrer et faire tourner}

\quoi
\textbf{ghcr.io} est le registre d'images de GitHub : l'endroit où les images
construites sont publiées. \textbf{docker compose} décrit une pile de
conteneurs dans un fichier et la démarre en une commande.

\pourquoi
Pour le registre : il est intégré au dépôt, l'authentification est celle de
GitHub Actions, et les images sont versionnées avec le commit qui les a
produites.

Pour compose : monter quatre machines virtuelles demande du temps et de la
mémoire. Pour développer, pour démontrer, et pour faire tourner les tests
d'intégration dans la chaîne, on veut la pile complète en une commande.

\role
Les images sont étiquetées trois fois : par version, par commit, et
\cmd{latest}. Le fichier \fichier{docker-compose.yml} décrit la base,
l'application, le Faucheur, et les six constructions d'images sous un profil
séparé.

Deux mécanismes y méritent d'être signalés.

\textbf{L'ancre YAML} : un bloc nommé une fois, repris par l'application et le
Faucheur. Il devient impossible qu'ils divergent sur une variable.

\textbf{Le refus de démarrer sans secret} : la syntaxe
\cmd{\$\{POSTGRES\_PASSWORD:?manquant\}} fait échouer le démarrage si la
variable n'est pas définie. L'erreur arrive tout de suite, pas trois semaines
plus tard.

\begin{piege}{Un volume par service là où il en fallait un partagé}
Le Faucheur \emph{écrit} les archives, l'application les \emph{sert} au
téléchargement. Avec un volume par service, l'utilisateur verrait dans son
tableau de bord une archive que le serveur ne trouverait pas. Le même volume
nommé doit être monté dans les deux.
\end{piege}

\gros[0.86]{zoom-archives.png}{Les archives servies par l'application, écrites
par le Faucheur : deux services, un seul volume.}

\apport
Le projet se livre et se lance en une commande, en développement comme en
intégration continue.

% -----------------------------------------------------------------------------
\section*{Ce que la plateforme sait faire après la couche 6}

Elle se vérifie elle-même. Onze contrôles à chaque envoi, dont six de sécurité,
et aucune livraison si l'un échoue. Elle a trouvé dix défauts réels que la
relecture n'avait pas vus.
